% Lesson 08: shared student/instructor board-supported handout.
% Structure: lesson_07_handout_body(1).tex.
% Main content: lewis_ch3_handwritten_notes.tex, source pages 1--3.
% Objective and reading: lesson_08.md.
% The alpha, xi, and collision-count equations are explicitly identified as
% review from the attached Lesson 07 handout, not additions to Lewis notes.
%
% This body has the same integration contract as Lesson 07: load the common
% handout preamble, then input this file. It includes begin/end{document}.
% The required template commands/environments and KSU colors are unchanged.
% All diagrams are embedded editable TikZ; there are no external assets.
%
% Source-handling decisions:
% - Sketches remain qualitative; original hand-traced curve coordinates retained.
% - The moderator table is reproduced without supplying missing conditions/units.
% - The "no intermediate spectrum reactors" and LEU-boundary statements remain
%   explicitly attributed quotations in discussion/verification prompts.
% - Unclear/cropped fragments (including "constituents >= 10 [cropped]" and
%   the missing word in the U-238 inelastic-scattering sentence) are not completed.
% - Solution-only reminders about reintroducing diffusion and unfinished
%   marginal questions have not been converted into new course assertions.
% - Numerical checkpoint answers are arithmetic consequences of the given table.

\providecommand{\SourceEq}[1]{\par\nobreak\hfill{\sffamily\scriptsize\color{KSUGray}#1}\par}
\usetikzlibrary{arrows.meta,decorations.pathreplacing}

\setlength{\abovedisplayskip}{4pt plus 1pt minus 1pt}
\setlength{\belowdisplayskip}{4pt plus 1pt minus 1pt}
\setlength{\abovedisplayshortskip}{3pt plus 1pt}
\setlength{\belowdisplayshortskip}{3pt plus 1pt}
\setlength{\jot}{2pt}
\tcbset{handout base/.append style={before skip=2.5pt,after skip=2.5pt}}

% The curve coordinates below are copied from the transcription. Only their
% size, color, labels, and surrounding layout have been adapted for columns.
\providecommand{\LessonEightPureFuelSketch}{%
\begin{tikzpicture}[x=.88cm,y=.76cm,>=Latex,font=\scriptsize,
                   line cap=round,line join=round]
  \path[use as bounding box] (-.68,-.65) rectangle (8.18,3.07);
  \draw[->,semithick] (0,-.06)--(0,2.67);
  \draw[->,semithick] (-.1,0)--(8.0,0);
  \node[above] at (0,2.67) {$\eta(E)$};
  \draw (-.10,1.26)--(.10,1.26);
  \node[left=2pt] at (0,1.26) {$\sim2$};
  \foreach \x in {2.25,5.00,7.40}{\draw (\x,-.06)--(\x,.06);}
  \node[below=4pt] at (2.25,0) {$1\unit{eV}$};
  \node[below=4pt] at (5.00,0) {$0.1\unit{MeV}$};
  \node[below=4pt] at (7.40,0) {$10\unit{MeV}$};
  \draw[KSUPurple,thick] (0,1.26)--(2.17,1.26)
    .. controls (2.20,1.28) and (2.21,1.06) .. (2.25,.87)
    .. controls (2.28,.77) and (2.32,1.23) .. (2.39,1.26)
    .. controls (2.45,1.15) and (2.46,.88) .. (2.55,.87)
    .. controls (2.63,.94) and (2.65,1.18) .. (2.74,1.27)
    .. controls (2.80,1.08) and (2.84,.78) .. (2.94,.88)
    .. controls (3.02,.98) and (3.05,1.20) .. (3.14,1.27)
    .. controls (3.22,1.02) and (3.25,.73) .. (3.35,.85)
    .. controls (3.44,.94) and (3.50,1.16) .. (3.57,1.27)
    .. controls (3.67,1.02) and (3.70,.78) .. (3.80,.87)
    .. controls (3.91,1.00) and (3.91,1.20) .. (4.01,1.26)
    .. controls (4.08,1.03) and (4.13,.86) .. (4.20,.93)
    .. controls (4.28,1.02) and (4.31,1.19) .. (4.40,1.27)
    .. controls (4.47,1.02) and (4.51,.94) .. (4.60,1.04)
    .. controls (4.75,1.18) and (4.82,1.25) .. (5.00,1.26)
    .. controls (6.15,1.28) and (6.92,1.57) .. (7.56,2.38);
  \draw[KSUPurple,thin,densely dashed] (4.83,1.18)--(4.95,.22)--(5.14,1.27);
  \draw[KSUOrange,semithick,decorate,
        decoration={brace,amplitude=3pt,raise=2pt}]
        (2.24,1.66)--(4.98,1.66)
        node[midway,above=7pt,font=\scriptsize\bfseries]{avoid!};
  \node[align=center,text=KSUGray,inner sep=1pt] (nu) at (6.35,2.65)
       {larger $\nu(E)$};
  \draw[->,KSUGray] (nu.south east) to[out=-25,in=150] (7.25,2.05);
\end{tikzpicture}%
}
\providecommand{\LessonEightMixtureSketch}{%
\begin{tikzpicture}[x=.85cm,y=.73cm,>=Latex,font=\scriptsize,
                   line cap=round,line join=round]
  \path[use as bounding box] (-.72,-1.16) rectangle (8.40,3.13);
  \draw[->,semithick] (0,-.06)--(0,2.65);
  \draw[->,semithick] (-.1,0)--(8.27,0);
  \node[above] at (0,2.65) {$\eta(E)$};
  \draw (-.10,1.40)--(.10,1.40);
  \node[left=2pt] at (0,1.40) {$\sim2$};
  \draw[KSUOrange] (-.10,.88)--(.10,.88);
  \node[left=2pt,text=KSUOrange] at (0,.88) {$\sim1.4$};
  \draw[KSUGray,densely dashed] (2.36,0)--(2.36,2.15);
  \draw[KSUGray,densely dashed] (6.04,0)--(6.04,2.15);
  \draw[KSUPurple,thick] (0,1.40)--(2.35,1.40)
    .. controls (2.45,1.17) and (2.48,.75) .. (2.60,.85)
    .. controls (2.70,1.04) and (2.75,1.29) .. (2.82,1.34)
    .. controls (2.91,1.13) and (2.94,.82) .. (3.05,.89)
    .. controls (3.15,1.02) and (3.20,1.30) .. (3.26,1.32)
    .. controls (3.37,.93) and (3.51,.65) .. (3.66,.80)
    .. controls (3.75,.97) and (3.82,1.25) .. (3.89,1.34)
    .. controls (3.96,1.12) and (3.96,.91) .. (4.08,.91)
    .. controls (4.21,.96) and (4.21,1.27) .. (4.30,1.32)
    .. controls (4.34,1.00) and (4.39,.79) .. (4.51,.94)
    .. controls (4.61,1.12) and (4.70,1.43) .. (4.80,1.38)
    .. controls (4.94,1.26) and (4.94,1.53) .. (5.02,1.41)
    .. controls (5.17,1.08) and (5.70,.99) .. (6.10,1.07)
    .. controls (7.11,1.26) and (7.44,1.54) .. (8.12,2.54);
  \draw[KSUOrange,thick] (0,.88)
    .. controls (.18,.86) and (.23,.97) .. (.34,.88)
    .. controls (.47,.81) and (.57,.91) .. (.71,.87)
    .. controls (.91,.83) and (.96,.91) .. (1.09,.86)
    .. controls (1.25,.80) and (1.36,.93) .. (1.49,.85)
    .. controls (1.65,.78) and (1.82,.91) .. (1.96,.86)
    --(2.28,.86)
    .. controls (2.36,.81) and (2.38,.27) .. (2.48,.24)
    .. controls (2.62,.22) and (2.76,.47) .. (2.83,.57)
    .. controls (2.87,.42) and (2.96,.35) .. (3.05,.44)
    .. controls (3.13,.30) and (3.29,.27) .. (3.37,.29)
    .. controls (3.48,.37) and (3.51,.28) .. (3.66,.27)
    .. controls (3.85,.18) and (3.88,.38) .. (4.01,.27)
    .. controls (4.17,.20) and (4.28,.27) .. (4.41,.22)
    .. controls (4.55,.21) and (4.59,.42) .. (4.72,.37)
    .. controls (4.81,.52) and (4.84,.25) .. (4.98,.27)
    .. controls (5.13,.23) and (5.20,.34) .. (5.34,.29)
    .. controls (5.57,.23) and (5.64,.29) .. (5.85,.27)
    .. controls (6.89,.30) and (7.49,.92) .. (7.87,1.63);
  \node[text=KSUPurple,inner sep=1pt] (twenty) at (6.64,2.72) {$20\%$};
  \draw[->,KSUPurple] (twenty.east) to[out=0,in=140] (8.01,2.31);
  \node[text=KSUOrange,anchor=east,inner sep=1pt] (nat) at (7.20,.58) {natural U};
  \draw[->,KSUOrange] (nat.east) to[out=0,in=-70] (7.64,1.08);
  \node[below=4pt] at (2.36,0) {$\sim1\unit{eV}$};
  \node[below=4pt] at (6.04,0) {$\sim0.1\unit{MeV}$};
  \node[below=4pt] at (8.18,0) {$E$};
  \node[text=KSUDeepPurple] at (1.0,-.91) {thermal};
  \node[text=KSUGray] at (4.1,-.91) {intermediate};
  \node[text=KSUDeepPurple] at (7.3,-.91) {fast};
\end{tikzpicture}%
}

\HandoutSetup
  {NE 630}
  {08}
  {Fuels, Moderators, and Reactor Spectra}
  {Lewis \S\S3.1--3.3}

\begin{document}
\MakeHandoutTitle

\textcolor{KSUDeepPurple}{\textbf{Primary objective.}}
Explain the different goals of thermal- and fast-spectrum reactor designs
using $\eta(E)$, $\xi$, $\Sigma_s$, and $\Sigma_a$.
\par\vspace{5pt}

% ===================== PAGE 1: ENERGY AND FUELS =====================
\begin{handoutcolumns}
\HandoutSection{1}{Energy and the multiplication factor}
Recall the generation-to-generation definition:
$$
k=\frac{\#\text{ neutrons in generation }i}
       {\#\text{ neutrons in generation }i-1}.
$$
What determines this value? Chapter 3 focuses on \textbf{energy}:
neutron cross sections are strongly energy dependent.

\begin{fixedbox}{The model for this chapter}
\renewcommand{\arraystretch}{1.10}
\begin{tabularx}{\linewidth}{@{}l >{\raggedright\arraybackslash}X@{}}
\textbf{Infinite} & Set aside spatial diffusion and leakage.\\
\textbf{Homogeneous} & Use a spatially uniform composition.\\
\textbf{Steady state} & Set aside time-dependent phenomena, including delayed-neutron dynamics.
\end{tabularx}
\medskip
The notes write the multiplication factor as $k_\infty$.
\end{fixedbox}

\textbf{Chapter roadmap.} Fuels $\to$ moderators $\to$ energy distributions
$\to$ energy averaging $\to$ $k_\infty$ in terms of
energy-averaged (``effective'') cross sections.

\textbf{Energy range of interest.} Roughly
$$
0.001\unit{eV}\ \lesssim E\ \lesssim\ 10\unit{MeV}.
$$

\HandoutSection{2}{Fuel: neutrons produced per absorption}
\begin{fixedbox}{The energy-dependent ratio $\eta(E)$}
$$
\eta(E)=\frac{\nu(E)\,\Sigma_f(E)}{\Sigma_a(E)}
       =\frac{\text{fission neutrons produced}}{\text{neutrons absorbed}}.
$$
\par
\renewcommand{\arraystretch}{1.10}
\begin{tabularx}{\linewidth}{@{}l >{\raggedright\arraybackslash}X@{}}
$\nu(E)$ & Mean neutrons produced per fission.\\
$\Sigma_f(E)$ & Macroscopic fission cross section.\\
$\Sigma_a(E)$ & Macroscopic absorption cross section.
\end{tabularx}\par
Use the same fuel composition in the numerator and denominator.
\end{fixedbox}

\begin{checkpoint}[What do the two ratios count?]
How does $\eta(E)$ differ from $k_\infty$? Which compares generations,
and which counts production per absorption at a specified energy?
\RuledLines{2}
\end{checkpoint}

\begin{boardbox}{Board work: absorption in a fuel mixture}
What nuclide densities and microscopic data are needed to construct
$\Sigma_a(E)$ for a uranium mixture?
\RuledLines{3}
\end{boardbox}

\switchcolumn
\RaggedRight
\HandoutSection{3}{Read the fuel sketches}
\begin{fixedbox}{Pure fissile material}
\centering
\LessonEightPureFuelSketch\par
\end{fixedbox}

\begin{fixedbox}{Natural uranium and an enriched mixture}
\centering
\LessonEightMixtureSketch\par
\end{fixedbox}
{\sffamily\scriptsize\itshape\color{KSUGray}
Note, these are cartoons to illustrate the concepts!\par}



\begin{checkpoint}[Ponderable]
If $\eta(E)$ represents an upper bound on $k_{\infty}$, why do you think
we rarely see reactors operating in the intermediate spectrum?
\RuledLines{2}
\end{checkpoint}

\HandoutSection{4}{Enrichment: How Much is Enough?}
Naval reactors typically use very high enrichment,
while civilian fuels range from natural uranium to $20\%$ enrichment,
the arbitrary(?) cutoff between low-enriched uranium \textbf{LEU} and \textbf{HEU}.
Traditional power reactors typically use enrichments below $5\%$; enrichments
between 5 and 20\% are known as High-Assay LEU or HALEU.

\begin{checkpoint}[Reading follow-up]
Why might a design favor high enrichment?
Why might a design favor low enrichment?
\RuledLines{2}
\end{checkpoint}
\end{handoutcolumns}

\newpage

% ==================== PAGE 2: MODERATION AND DESIGN =================
\begin{handoutcolumns}
\HandoutSection{5}{Thermal and fast design strategies}
\begin{fixedbox}{Thermal-spectrum reactors}
Use \textbf{low-mass moderators} to move neutrons from fast to thermal
energies with as \textbf{few collisions in the resonance range} as possible.
The notes connect suitable fuel--moderator choices with low or no enrichment.
\end{fixedbox}

\begin{fixedbox}{Fast-spectrum reactors}
Choose materials that do not slow neutrons effectively:
\textbf{avoid low-mass materials} and \textbf{minimize inelastic scattering}.
\end{fixedbox}

\textbf{Keep the mechanisms distinct.} The moderator measures below concern
\textbf{elastic} scattering. At higher energies, \textbf{inelastic}
scattering becomes significant, affecting fast spectra.

\HandoutSection{6}{What makes a good moderator?}
\begin{fixedbox}{Three desirable properties}
\renewcommand{\arraystretch}{1.16}
\begin{tabularx}{\linewidth}{@{}l >{\raggedright\arraybackslash}X@{}}
\textbf{Large $\xi$} & A light target gives a large average logarithmic energy loss per elastic collision.\\
\textbf{Large $\Sigma_s$} & Large number density and/or large $\sigma_s$ favors collisions in the moderator.\\
\textbf{Small $\Sigma_a$} & Prefer scattering to absorption; the notes emphasize $\sigma_a\ll\sigma_s$.
\end{tabularx}
\end{fixedbox}

\begin{fixedbox}{Review from Lesson 07: decrement and collision count}
For elastic scattering isotropic in the CM frame, with target-to-neutron
mass ratio $A$,
$$
\alpha=\frac{(A-1)^2}{(A+1)^2},\qquad
\xi=1+\frac{\alpha\ln\alpha}{1-\alpha}.
$$
At $A=1$, use the limiting value $\xi=1$.
$$
n\approx\frac{1}{\xi}\ln\frac{E_0}{E_n}.
$$
For compounds, use the effective mixture decrement; the Lesson 07
water example uses $\xi_{\mathrm{water}}=0.93$.
\end{fixedbox}

\begin{checkpoint}[One favorable property is not enough]
Which moderator in the table gives the fewest collisions for a specified
$E_0/E_n$? Does the same moderator have the largest slowing-down ratio?
\RuledLines{2}
\end{checkpoint}

\begin{boardbox}{Notes: connecting the three properties}
\RuledLines{3}
\end{boardbox}

\switchcolumn
\RaggedRight
\HandoutSection{7}{Slowing-down power and ratio}
\begin{fixedbox}{Combine the moderator properties}
$$
\text{slowing-down power}=\xi\Sigma_s,
$$
$$
\text{slowing-down ratio}=\frac{\xi\Sigma_s}{\Sigma_a}.
$$
The first combines energy loss and scattering; the second compares that
combination with absorption.
\end{fixedbox}

\begin{fixedbox}{Moderator comparison}
\centering
\renewcommand{\arraystretch}{1.20}
\setlength{\tabcolsep}{5pt}
\begin{tabular}{@{}lrrr@{}}
\hline
 & $\xi$ & $\xi\Sigma_s$ & $\xi\Sigma_s/\Sigma_a$\\
\hline
$\mathrm{H_2O}$ & $0.93$ & $1.28$ & \textcolor{KSUOrange}{\textbf{58}}\\
$\mathrm{D_2O}$ & $0.51$ & $0.18$ & \textcolor{KSUDeepPurple}{\textbf{21,000}}\\
$\mathrm{C}$ & $0.158$ & $0.056$ & \textcolor{KSUDeepPurple}{\textbf{200}}\\
\hline
\end{tabular}\par
\medskip
\RaggedRight
{\footnotesize\textit{Source limitation.} Values are reproduced as transcribed.
The notes do not specify reference conditions or the units used for
$\xi\Sigma_s$.}\par
\end{fixedbox}

\textbf{Natural-uranium connection.} Heavy water and graphite
are moderators with which natural-U fuel can be used; graphite
was using in Fermi's pile.

\HandoutSection{8}{Use the measures}
\begin{checkpoint}[Recover absorption from the table]
Rearrange the slowing-down ratio to obtain
$$
\Sigma_a=\MathBlank[1.25in]{\frac{\xi\Sigma_s}{\xi\Sigma_s/\Sigma_a}}.
$$
Then compute the implied absorption cross sections:
\begin{align*}
\mathrm{H_2O}:\quad&\Sigma_a=\frac{1.28}{58}
                 \approx\MathBlank[.85in]{2.21\times10^{-2}},\\[5pt]
\mathrm{D_2O}:\quad&\Sigma_a=\frac{0.18}{21{,}000}
                 \approx\MathBlank[.85in]{8.57\times10^{-6}}.
\end{align*}
The results have the same inverse-length units as the table's
slowing-down power; no numerical unit convention is supplied here.
\end{checkpoint}

\begin{checkpoint}[Compare the design choices]
Why can water have the largest slowing-down power in the table but the
smallest slowing-down ratio? How does this connect to different fuel
choices?
\RuledLines{3}
\end{checkpoint}

\end{handoutcolumns}


\end{document}
